\begin{proof}[Proof of \cref{obj:lem:upper-risk}]
\begin{enumerate}
\item Define the numerical constants
\[
\begin{aligned}
C_{\mathrm{bias}}&=268435456,\\
C_{\mathrm{switch}}&=201719808,\\
C_{\mathrm{log}}&=720(4/3)^6(e+1),\\
C_{\mathrm{non}}&=
8\bigl(C_{\mathrm{switch}}^2C_{\mathrm{log}}+119\bigr)
+2C_{\mathrm{bias}}^2+104,\\
C_{\mathrm{fall}}&=4e^{128},\\
C_U&=C_{\mathrm{non}}+C_{\mathrm{fall}}.
\end{aligned}
\]
These constants are finite and positive. Fix parameters and a law satisfying the hypotheses, and put
\[
\delta=1-q\in[0,1/2],
\qquad
\ell_n=\log(e+n).
\]
By the independent sampling hypothesis,
\[
(O_1,\ldots,O_n)\sim P_O^{\otimes n}.
\]
From \cref{obj:def:rate},
\[
r_{n,d,q}\ge \frac1n>0.
\]
We establish the risk bounds \(C_{\mathrm{fall}}r_{n,d,q}\) and
\(C_{\mathrm{non}}r_{n,d,q}\) in the corresponding branches of
\cref{obj:def:mm-estimator}. Each is at most \(C_Ur_{n,d,q}\).
% lean: CausalSmith.Stat.MarNearcompleteFrontier.upper_risk

\item Suppose first that \(q=1\). The arrival floor makes the conditional arrival probability equal to one in every positive-probability cell. Consequently,
\[
P(R=0)=0.
\]
Define the complete-arrival score on the observed-record space by
\[
F_{\mathrm{complete}}:O\longrightarrow[-2,2],
\qquad
F_{\mathrm{complete}}(O)=2(2A-1)RY.
\]
The restrictions defining \cref{obj:def:law-class} give
\[
\begin{aligned}
\mathbb E_P[F_{\mathrm{complete}}(O)]
&=2P(A=1,Y=1)-2P(A=0,Y=1)\\
&=P(Y(1)=1)-P(Y(0)=1)
=\tau(P).
\end{aligned}
\]
Indeed, consistency replaces the realized outcome within each arm by its potential outcome, and randomized independence and balanced assignment give
\(2P(A=a,Y=1)=P(Y(a)=1)\) for \(a\in\{0,1\}\).
Also,
\[
\operatorname{Var}_P(F_{\mathrm{complete}}(O))
\le \mathbb E_P[F_{\mathrm{complete}}(O)^2]\le4.
\]
Independence of the observations therefore yields
\[
\mathbb E_P\!\left[
\left\{\frac1n\sum_{i=1}^nF_{\mathrm{complete}}(O_i)-\tau(P)\right\}^2
\right]
=\frac{\operatorname{Var}_P(F_{\mathrm{complete}}(O))}{n}
\le\frac4n.
\]
The binary potential outcomes imply \(\tau(P)\in[-1,1]\), and interval projection satisfies
\[
\left|\Pi_{[-1,1]}(t)-\tau(P)\right|
\le |t-\tau(P)|,
\qquad t\in\mathbb R.
\]
Thus the complete-arrival branch satisfies
\[
\mathbb E_P[(\widehat\tau_{\mathrm{MM}}-\tau(P))^2]
\le \frac4n
=4r_{n,d,1}
\le C_{\mathrm{fall}}r_{n,d,1}
\le C_Ur_{n,d,1}.
\]
% lean: CausalSmith.Stat.MarNearcompleteFrontier.upper_complete_arrival_risk_rate

\item Suppose next that
\[
q\ne1,
\qquad
\ell_n<128\ \text{or}\ d\ge n\ell_n.
\]
When \(\ell_n<128\), exponentiating its definition gives
\[
n<e+n=e^{\ell_n}<e^{128}.
\]
Both the estimator and the target lie in \([-1,1]\), so
\[
\mathbb E_P[(\widehat\tau_{\mathrm{MM}}-\tau(P))^2]
\le4
\le\frac{4e^{128}}n
\le C_{\mathrm{fall}}r_{n,d,q}.
\]

For the large-dimension branch, introduce the following quantities, defined for every cell, including null cells:
\[
\begin{gathered}
\mathcal J=\{0,\ldots,d-1\}\times\{0,1\}^2,
\qquad j=(x,a,s)\in\mathcal J,\\
\pi_j=P(X=x,A=a,S=s),\\
t_j=P(X=x,A=a,S=s,R=1),
\qquad
w_j=\pi_j-t_j=P(X=x,A=a,S=s,R=0),\\
\mu_P(j)=\mu_P(x,a,s),
\qquad
\eta_j=\mu_P(j)-\frac12,
\qquad
\epsilon_j=2a-1.
\end{gathered}
\]
The mean \(\mu_P(j)\) uses the zero-denominator convention in
\cref{obj:def:identified-functional}. In particular,
\[
|\eta_j|\le\frac12,\qquad
\sum_{j\in\mathcal J}\pi_j=1.
\]
The arrival floor gives \(q\pi_j\le t_j\) on positive cells. On a null cell,
\(\pi_j=t_j=w_j=0\), so the same inequality holds. Since \(q\ge1/2\),
\[
0\le w_j\le\delta\pi_j\le2\delta t_j,
\qquad
\sum_{j\in\mathcal J}w_j\le\delta.
\]
Balanced assignment gives
\[
\sum_{j\in\mathcal J}\epsilon_j\pi_j=0.
\]
Applying \cref{obj:lem:identification-infrastructure} and centering the cell means therefore gives
\[
\tau(P)
=2\sum_{j\in\mathcal J}\epsilon_j\pi_j\eta_j
=2\sum_{j\in\mathcal J}\epsilon_j(t_j\eta_j+w_j\eta_j).
\]

Define the fallback score by
\[
F_{\mathrm{fallback}}:O\longrightarrow[-1,1],
\qquad
F_{\mathrm{fallback}}(O)=2(2A-1)R(RY-1/2).
\]
On the support of \(P_O\), its mean is
\[
\mathbb E_P[F_{\mathrm{fallback}}(O)]
=2\sum_{j\in\mathcal J}\epsilon_jt_j\eta_j.
\]
Hence
\[
\left|\mathbb E_P[F_{\mathrm{fallback}}(O)]-\tau(P)\right|
=\left|2\sum_{j\in\mathcal J}\epsilon_jw_j\eta_j\right|
\le\sum_{j\in\mathcal J}w_j
\le\delta.
\]
The score bound and independent sampling imply
\[
\mathbb E_P\!\left[
\left\{\frac1n\sum_{i=1}^nF_{\mathrm{fallback}}(O_i)
-\mathbb E_P[F_{\mathrm{fallback}}(O)]\right\}^2
\right]\le\frac1n.
\]
Using projection contraction and
\((u+v)^2\le2u^2+2v^2\), we obtain
\[
\mathbb E_P[(\widehat\tau_{\mathrm{MM}}-\tau(P))^2]
\le\frac2n+2\delta^2.
\]
For \(d\ge n\ell_n\), the minimum in \cref{obj:def:rate} equals one. Thus
\[
\frac2n+2\delta^2
=2r_{n,d,q}
\le C_{\mathrm{fall}}r_{n,d,q}.
\]
Both fallback subcases consequently satisfy the bound with \(C_U\).
The cell quantities just defined will also be used below.
% lean: CausalSmith.Stat.MarNearcompleteFrontier.upper_fallback_branches_risk_rate

\item Consider now the remaining regime
\[
q\ne1,\qquad \ell_n\ge128,\qquad d<n\ell_n.
\]
The tuning parameters and benchmark become
\[
m=\frac n8,\qquad
k=\lfloor\ell_n/64\rfloor,\qquad
B=256\ell_n,
\]
\[
\frac{\ell_n}{128}\le k\le\frac{\ell_n}{64}
=\frac{B}{16384},
\qquad
g_{n,d,q}=\frac{\delta d}{n\ell_n},
\qquad
r_{n,d,q}=\frac1n+g_{n,d,q}^2.
\]
The lower bound on \(k\) follows from
\(\lfloor\ell_n/64\rfloor\ge\ell_n/64-1\ge\ell_n/128\).

Extend the observed sample to an infinite independent sequence with law \(P_O\). Independently draw
\[
N\sim\operatorname{Poisson}(n/2),
\qquad
Z_i\stackrel{\mathrm{iid}}{\sim}\operatorname{Unif}\{1,2,3,4\},
\qquad i\ge1,
\]
and define the uncapped stream sets
\[
\mathcal I_t^\star=\{i\le N:Z_i=t\},
\qquad t\in\{1,2,3,4\}.
\]
Write \(\mathbb E_\star\) and \(\Pr_\star\) for expectation and probability under this joint experiment. For an observed record \(o\), define its stream histogram count by
\[
K_{t,o}^{\mathrm{hist}}
=\sum_{i\in\mathcal I_t^\star}\mathbf1\{O_i=o\}.
\]
For \(u_{t,o}\in[0,1]\), conditioning on \(N\) gives the joint probability generating function
\[
\mathbb E_\star\!\left[
\prod_{t=1}^4\prod_{o\in O}
u_{t,o}^{K_{t,o}^{\mathrm{hist}}}\right]
=
\exp\!\left\{
m\sum_{t=1}^4\sum_{o\in O}P_O(\{o\})(u_{t,o}-1)
\right\}.
\]
Its factorization shows that these histogram counts are independent Poisson variables with means \(mP_O(\{o\})\).

Evaluate the correction formulas of \cref{obj:def:mm-estimator} on the uncapped streams:
\[
\begin{aligned}
V_j^\star&=\frac1m\sum_{i\in\mathcal I_2^\star}
\mathbf1\{(X_i,A_i,S_i)=j,\ R_i=0\},\\
C_j^{\prime\star}&=\sum_{i\in\mathcal I_3^\star}
\mathbf1\{(X_i,A_i,S_i)=j,\ R_i=1\},\\
C_j^\star&=\sum_{i\in\mathcal I_4^\star}
\mathbf1\{(X_i,A_i,S_i)=j,\ R_i=1\},\\
U_j^\star&=\sum_{i\in\mathcal I_4^\star}
\mathbf1\{(X_i,A_i,S_i)=j,\ Y_i^{\mathrm{obs}}=1\},\\
H_j^\star&=\sum_{v=1}^{k-1}
a_v(U_j^\star-C_j^\star/2)(C_j^\star-1)_{v-1},\\
D_j^\star&=
\begin{cases}
U_j^\star/C_j^\star-1/2,&C_j^\star>0,\\
0,&C_j^\star=0,
\end{cases}\\
G_j^\star&=
H_j^\star\mathbf1\{C_j^{\prime\star}\le B/4\}
+D_j^\star\mathbf1\{C_j^{\prime\star}>B/4\}.
\end{aligned}
\]
Define the linear term, aggregate correction, raw estimate, and projected ideal estimate by
\[
\begin{aligned}
\tau_{\mathrm{lin}}
&=\frac2m\sum_{i\in\mathcal I_1^\star}
(2A_i-1)R_i(Y_i^{\mathrm{obs}}-1/2),\\
W_{\mathrm{corr}}
&=\sum_{j\in\mathcal J}\epsilon_jV_j^\star G_j^\star,\\
\tau_{\mathrm{raw}}
&=\tau_{\mathrm{lin}}+2W_{\mathrm{corr}},\\
\widetilde\tau_\star&=\Pi_{[-1,1]}(\tau_{\mathrm{raw}}).
\end{aligned}
\]

The capped randomized estimate \(\widetilde\tau\) in
\cref{obj:def:mm-estimator} agrees with \(\widetilde\tau_\star\) when \(N\le n\).
On the complementary event their squared errors are bounded by four.
Conditional Jensen and this coupling therefore give
\[
\begin{aligned}
\mathbb E_P[(\widehat\tau_{\mathrm{MM}}-\tau(P))^2]
&\le\mathbb E_\star[(\widetilde\tau-\tau(P))^2]\\
&\le\mathbb E_\star[(\widetilde\tau_\star-\tau(P))^2]
+4\Pr_\star(N>n).
\end{aligned}
\]
Exponential Markov, applied to \(2^N\), yields
\[
\Pr_\star(N>n)
\le2^{-n}\mathbb E_\star[2^N]
=\exp\{-(\log2-1/2)n\}
\le e^{-n/8}
\le\frac8n.
\]
Here \(\log2-1/2\ge1/8\), and \(n/8\le e^{n/8}\) gives the last inequality. Consequently,
\[
\mathbb E_P[(\widehat\tau_{\mathrm{MM}}-\tau(P))^2]
\le\mathbb E_\star[(\widetilde\tau_\star-\tau(P))^2]
+\frac{32}{n}.
\]
% lean: CausalSmith.Stat.MarNearcompleteFrontier.upper_nonfallback_risk_le_ideal_add_rate

\item All the correction statistics have finite second moments: their polynomial parts are finite polynomials in Poisson counts, and their ratio parts are bounded. Projection contraction and the variance--bias identity give
\[
\mathbb E_\star[(\widetilde\tau_\star-\tau(P))^2]
\le
\operatorname{Var}_\star(\tau_{\mathrm{raw}})
+\{\mathbb E_\star[\tau_{\mathrm{raw}}]-\tau(P)\}^2.
\]
To bound the first term, define
\[
\begin{aligned}
K_{\mathrm{lin},+}
&=\sum_{i\in\mathcal I_1^\star}
\mathbf1\{R_i=1,\ A_i=Y_i^{\mathrm{obs}}\},\\
K_{\mathrm{lin},-}
&=\sum_{i\in\mathcal I_1^\star}
\mathbf1\{R_i=1,\ A_i\ne Y_i^{\mathrm{obs}}\}.
\end{aligned}
\]
Then
\[
\tau_{\mathrm{lin}}
=\frac{K_{\mathrm{lin},+}-K_{\mathrm{lin},-}}m.
\]
Each count is Poisson with mean at most \(m\). Using
\(\operatorname{Var}(U-V)\le2\operatorname{Var}(U)+2\operatorname{Var}(V)\),
\[
\operatorname{Var}_\star(\tau_{\mathrm{lin}})
\le\frac{4m}{m^2}
=\frac{32}{n}.
\]
Applying the corresponding sum inequality to
\(\tau_{\mathrm{raw}}=\tau_{\mathrm{lin}}+2W_{\mathrm{corr}}\) gives
\[
\operatorname{Var}_\star(\tau_{\mathrm{raw}})
\le\frac{64}{n}
+8\operatorname{Var}_\star(W_{\mathrm{corr}}).
\]
% lean: CausalSmith.Stat.MarNearcompleteFrontier.upper_fourStreamRawEstimate_variance_le_aggregate

\item For every \(j\in\mathcal J\), define its arrived-count intensity and light-branch probability by
\[
\zeta_j=mt_j\ge0,
\qquad
\wp_j=\Pr_\star(C_j^{\prime\star}\le B/4)\in[0,1].
\]
Poisson splitting gives
\[
\mathbb E_\star[V_j^\star]=w_j,
\qquad
\mathbb E_\star[(V_j^\star)^2]=w_j^2+\frac{w_j}{m},
\]
\[
C_j^{\prime\star},C_j^\star\sim\operatorname{Poisson}(\zeta_j),
\qquad
U_j^\star\mid C_j^\star
\sim\operatorname{Binomial}(C_j^\star,\mu_P(j)).
\]
The conditional binomial statement includes \(t_j=0\): then both counts vanish almost surely and \(\mu_P(j)=0\).

For integers \(v\ge1\), use the falling factorial
\[
(C_j^\star)_v=\prod_{t=0}^{v-1}(C_j^\star-t).
\]
Conditional expectation of \(U_j^\star\), followed by the Poisson factorial-moment identity, gives
\[
\begin{aligned}
\mathbb E_\star[
(U_j^\star-C_j^\star/2)(C_j^\star-1)_{v-1}]
&=\eta_j\mathbb E_\star[(C_j^\star)_v]\\
&=\eta_j\zeta_j^v.
\end{aligned}
\]
Thus the coefficient identity in \cref{obj:def:mm-estimator} yields
\[
\mathbb E_\star[H_j^\star]
=\eta_j\{1-Q_k(\zeta_j)\}.
\]
Similarly, conditioning on \(C_j^\star>0\) gives
\[
\mathbb E_\star[D_j^\star]
=\eta_j\Pr_\star(C_j^\star>0)
=\eta_j(1-e^{-\zeta_j}).
\]
Independence of the pilot stream from the fourth stream now implies
\[
\mathbb E_\star[G_j^\star]-\eta_j
=-\eta_j\{\wp_jQ_k(\zeta_j)+(1-\wp_j)e^{-\zeta_j}\}.
\]
The first-stream mean and independence of the second stream from the pilot and fourth streams also give
\[
\mathbb E_\star[\tau_{\mathrm{lin}}]
=2\sum_{j\in\mathcal J}\epsilon_jt_j\eta_j,
\qquad
\mathbb E_\star[V_j^\star G_j^\star]
=w_j\mathbb E_\star[G_j^\star].
\]
Combining these identities with the cell decomposition of the target,
\[
\mathbb E_\star[\tau_{\mathrm{raw}}]-\tau(P)
=
2\sum_{j\in\mathcal J}
\epsilon_jw_j\{\mathbb E_\star[G_j^\star]-\eta_j\}.
\]
% lean: CausalSmith.Stat.MarNearcompleteFrontier.upper_fourStreamRawEstimate_bias_eq

\item We next derive the polynomial and pilot bounds used to control this bias. For \(0<z\le B\), the argument \(1-2z/B\) belongs to \([-1,1]\), so the Chebyshev bound \(|T_k(1-2z/B)|\le1\) gives
\[
Q_k(z)\ge0,
\qquad
zQ_k(z)
=\frac{B\{1-T_k(1-2z/B)\}}{2k^2}
\le\frac{B}{k^2}.
\]
At \(z=0\), the convention is \(Q_k(0)=1\).
The degree bound already established implies
\[
\frac{B}{k^2}
\le\frac{4194304}{\ell_n}.
\]

For the second-moment estimates, define the coefficient norm of a polynomial
\(\mathcal P(u)\) by
\[
\|\mathcal P\|_{\mathrm{coeff},1}
=\sum_{v\ge0}|[u^v]\mathcal P|,
\]
where \([u^v]\mathcal P\) denotes its coefficient of \(u^v\); the sum is finite. Define the shifted Chebyshev polynomials by
\[
\mathcal P_v(u)=T_v(1-2u),
\qquad v\in\mathbb N_0.
\]
Their recurrence and initial values give
\[
\mathcal P_{v+1}(u)
=2(1-2u)\mathcal P_v(u)-\mathcal P_{v-1}(u)
\quad(v\ge1),
\]
\[
\|\mathcal P_0\|_{\mathrm{coeff},1}=1,
\qquad
\|\mathcal P_1\|_{\mathrm{coeff},1}=3,
\]
\[
\|\mathcal P_{v+1}\|_{\mathrm{coeff},1}
\le6\|\mathcal P_v\|_{\mathrm{coeff},1}
+\|\mathcal P_{v-1}\|_{\mathrm{coeff},1}.
\]
Induction therefore yields
\[
\|\mathcal P_v\|_{\mathrm{coeff},1}\le7^v.
\]
Indeed, the induction step uses
\(6\cdot7^v+7^{v-1}\le7^{v+1}\).
Since
\[
Q_k(Bu)=\frac{1-\mathcal P_k(u)}{2k^2u},
\]
with the removable value at zero, deleting the zero constant term and dividing by \(u\) cannot increase the coefficient norm. Hence
\[
\sum_{v=1}^{k-1}|a_v|B^v
\le\frac{1+7^k}{2k^2}
\le7^k
\le8^k.
\]

On observed-law support, \(0\le U_j^\star\le C_j^\star\), and
\[
|(U_j^\star-C_j^\star/2)(C_j^\star-1)_{v-1}|
\le\frac{(C_j^\star)_v}{2}
\le\frac{(C_j^\star)^v}{2}.
\]
For every \(z\ge0\) and \(1\le v\le k\),
\[
(z/B)^v\le\max\{1,(z/B)^k\}.
\]
The coefficient bound therefore gives the pointwise envelope
\[
(H_j^\star)^2
\le\frac{64^k}{4}
\max\left\{1,\left(\frac{C_j^\star}{B}\right)^{2k}\right\}
\le\frac{64^k}{4}
\left\{1+\left(\frac{C_j^\star}{B}\right)^{2k}\right\}.
\]
For an integer \(v\ge0\), expansion of an ordinary Poisson moment into factorial moments gives
\[
\mathbb E_\star[(C_j^\star)^v]
=\sum_{t=0}^v {v\brace t}\zeta_j^t
\le\sum_{t=0}^v\binom vt v^{\,v-t}\zeta_j^t
=(\zeta_j+v)^v.
\]
Here \({v\brace t}\) counts partitions into \(t\) nonempty blocks: choosing one representative per block and assigning the remaining elements gives
\({v\brace t}\le\binom vt v^{v-t}\). The case \(v=0\) uses the constant zeroth moment. Integrating the pointwise envelope yields
\[
\mathbb E_\star[(H_j^\star)^2]
\le\frac{64^k}{4}
\left\{1+\left(\frac{\zeta_j+2k}{B}\right)^{2k}\right\}
\le\frac{64^k}{2}
\max\left\{1,\left(\frac{\zeta_j+2k}{B}\right)^{2k}\right\}.
\]
% lean: CausalSmith.Stat.MarNearcompleteFrontier.upper_streamH_second_moment_le

\item When \(\zeta_j\le B\), the degree calibration gives
\[
0\le\frac{\zeta_j+2k}{B}\le\frac32\le e^{1/2}.
\]
Consequently,
\[
\mathbb E_\star[(H_j^\star)^2]
\le\frac{64^k}{2}e^k
\le e^{7k}
\le e^{\ell_n/8},
\]
where \(64^k\le e^{6k}\) and \(7k\le\ell_n/8\).

For \(\zeta_j\ge B\), define
\[
\upsilon_j=\frac{\zeta_j+2k}{B}.
\]
From \(k\le B/16384\) and \(B\le\zeta_j\),
\[
2k\le\zeta_j,\qquad
k\upsilon_j\le\frac{\zeta_j}{8192},
\qquad
6k+2k\upsilon_j\le\frac{\zeta_j}{8}.
\]
Using \(\upsilon_j\le e^{\upsilon_j}\) in the second-moment envelope gives
\[
\mathbb E_\star[(H_j^\star)^2]
\le
64^k\max\{1,\upsilon_j^{2k}\}
\le e^{6k+2k\upsilon_j}
\le e^{\zeta_j/8}.
\]

Exponential Markov for the Poisson pilot count, using \(4^{-C_j^{\prime\star}}\) and \(4^{C_j^{\prime\star}}\), respectively, gives
\[
\wp_j\le
\exp\{B\log4/4-3\zeta_j/4\},
\qquad
1-\wp_j\le
\exp\{3\zeta_j-B\log4/4\}.
\]
Since \(1\le\log4\le2\), for \(\zeta_j\ge B\) these imply
\[
\wp_j\le e^{-\zeta_j/4},
\qquad
\wp_j\mathbb E_\star[(H_j^\star)^2]
\le e^{-\zeta_j/8}.
\]
Also, for every \(\zeta_j\ge0\),
\[
(1-\wp_j)e^{-\zeta_j}\le e^{-B/16}.
\]
To see this, if \(\zeta_j\ge B/16\), use \(1-\wp_j\le1\). Otherwise, the upper-tail bound gives
\[
(1-\wp_j)e^{-\zeta_j}
\le e^{2\zeta_j-B\log4/4}
\le e^{-B/16}.
\]

For \(\zeta_j\ge B\), Cauchy--Schwarz and \(\wp_j^2\le\wp_j\) yield
\[
\{\wp_j|\mathbb E_\star[H_j^\star]|\}^2
\le\wp_j\mathbb E_\star[(H_j^\star)^2]
\le e^{-\zeta_j/8},
\]
and hence
\[
\wp_j|\mathbb E_\star[H_j^\star]|
\le e^{-\zeta_j/16}.
\]
Using
\(\eta_jQ_k(\zeta_j)=\eta_j-\mathbb E_\star[H_j^\star]\), we obtain
\[
\begin{aligned}
\wp_j|\eta_jQ_k(\zeta_j)|
&\le\frac{\wp_j}{2}
+\wp_j|\mathbb E_\star[H_j^\star]|\\
&\le\frac32e^{-B/16}.
\end{aligned}
\]
The pilot bias identity therefore gives, on these large-intensity cells,
\[
|\mathbb E_\star[G_j^\star]-\eta_j|
\le2e^{-B/16}.
\]
% lean: CausalSmith.Stat.MarNearcompleteFrontier.upper_stream_pilot_large_bias_abs_le

\item We can now sum the bias. If \(\zeta_j=0\), then \(t_j=0\), and
\(0\le w_j\le2\delta t_j\) forces \(w_j=0\). Thus null arrived cells contribute zero to the weighted bias.

For \(0<\zeta_j\le B\), polynomial nonnegativity and
\(w_j\le2\delta t_j=2\delta\zeta_j/m\) give
\[
w_j\wp_j|Q_k(\zeta_j)|
\le w_jQ_k(\zeta_j)
\le\frac{2\delta}{m}\frac{B}{k^2}
\le\frac{2\delta}{m}\frac{4194304}{\ell_n}.
\]
There are \(4d\) cells in \(\mathcal J\), so
\[
\sum_{\substack{j\in\mathcal J\\\zeta_j\le B}}
w_j\wp_j|Q_k(\zeta_j)|
\le
4d\,\frac{2\delta}{m}\frac{4194304}{\ell_n}
=C_{\mathrm{bias}}g_{n,d,q}.
\]
For a small-intensity cell, the pilot bias identity and
\(|\eta_j|\le1/2\) give
\[
2w_j|\mathbb E_\star[G_j^\star]-\eta_j|
\le w_j\wp_j|Q_k(\zeta_j)|+w_je^{-B/16}.
\]
For a large-intensity cell, the preceding bound gives
\[
2w_j|\mathbb E_\star[G_j^\star]-\eta_j|
\le4w_je^{-B/16}.
\]
Together, these yield the uniform cellwise bound
\[
2w_j|\mathbb E_\star[G_j^\star]-\eta_j|
\le
\mathbf1\{\zeta_j\le B\}w_j\wp_j|Q_k(\zeta_j)|
+4w_je^{-B/16}.
\]
The triangle inequality and \(\sum_jw_j\le\delta\) therefore imply
\[
|\mathbb E_\star[\tau_{\mathrm{raw}}]-\tau(P)|
\le C_{\mathrm{bias}}g_{n,d,q}
+4\delta e^{-16\ell_n}.
\]
The remainder satisfies
\[
(4\delta e^{-16\ell_n})^2
=16\delta^2e^{-32\ell_n}
\le4e^{-32\ell_n}
\le4e^{-\ell_n}
=\frac4{e+n}
\le\frac4n.
\]
Squaring the bias bound with \((u+v)^2\le2u^2+2v^2\) gives
\[
\{\mathbb E_\star[\tau_{\mathrm{raw}}]-\tau(P)\}^2
\le
2C_{\mathrm{bias}}^2g_{n,d,q}^2+\frac8n
\le(2C_{\mathrm{bias}}^2+8)r_{n,d,q}.
\]
Combining projection, the raw variance bound, and the capping penalty established above,
\[
\begin{aligned}
\mathbb E_P[(\widehat\tau_{\mathrm{MM}}-\tau(P))^2]
&\le
8\operatorname{Var}_\star(W_{\mathrm{corr}})
+\frac{96}{n}
+(2C_{\mathrm{bias}}^2+8)r_{n,d,q}\\
&\le
8\operatorname{Var}_\star(W_{\mathrm{corr}})
+(2C_{\mathrm{bias}}^2+104)r_{n,d,q}.
\end{aligned}
\]
% lean: CausalSmith.Stat.MarNearcompleteFrontier.upper_nonfallback_risk_le_aggregate

\item The histogram coordinates used by two distinct cells are disjoint. Their independence, together with \(\epsilon_j^2=1\), gives
\[
\operatorname{Var}_\star(W_{\mathrm{corr}})
=\sum_{j\in\mathcal J}
\operatorname{Var}_\star(V_j^\star G_j^\star).
\]
We first bound the ratio component.

Since \(|D_j^\star|\le1/2\), its variance is at most \(1/4\). Thus, when \(\zeta_j\le15\),
\[
\operatorname{Var}_\star(D_j^\star)
\le\frac14\le\frac4{1+\zeta_j}.
\]
For \(\zeta_j>15\), define the arrived-success and arrived-failure fractions by
\[
F_{j,+}=
\begin{cases}
U_j^\star/C_j^\star,&C_j^\star>0,\\
0,&C_j^\star=0,
\end{cases}
\qquad
F_{j,-}=
\begin{cases}
(C_j^\star-U_j^\star)/C_j^\star,&C_j^\star>0,\\
0,&C_j^\star=0.
\end{cases}
\]
Then
\[
D_j^\star=\frac{F_{j,+}-F_{j,-}}2.
\]
Conditioning on the arrived count gives
\[
\begin{aligned}
\mathbb E_\star[(F_{j,+}-\mu_P(j))^2]
&\le
\Pr_\star(C_j^\star=0)
+\mathbb E_\star\!\left[
\frac{\mathbf1\{C_j^\star>0\}}{C_j^\star}\right],\\
\mathbb E_\star[(F_{j,-}-(1-\mu_P(j)))^2]
&\le
\Pr_\star(C_j^\star=0)
+\mathbb E_\star\!\left[
\frac{\mathbf1\{C_j^\star>0\}}{C_j^\star}\right].
\end{aligned}
\]
On a positive count, the conditional binomial variance divided by the squared count is at most its reciprocal; on a zero count, the squared error is at most one.

For a Poisson count of intensity \(\zeta_j\), the zero-safe reciprocal on the right is at most one. If \(\zeta_j\le1\), this is at most \(4/(1+\zeta_j)\). If \(\zeta_j>1\), the pointwise bounds
\[
\mathbf1\{C_j^\star=0\}
+\frac{\mathbf1\{C_j^\star>0\}}{C_j^\star}
\le\frac2{C_j^\star+1}
\]
and the Poisson identity
\[
\mathbb E_\star\!\left[\frac1{C_j^\star+1}\right]
=\frac{1-e^{-\zeta_j}}{\zeta_j}
\]
give
\[
\Pr_\star(C_j^\star=0)
+\mathbb E_\star\!\left[
\frac{\mathbf1\{C_j^\star>0\}}{C_j^\star}\right]
\le\frac2{\zeta_j}
\le\frac4{1+\zeta_j}.
\]
The reciprocal identity has value one at \(\zeta_j=0\), by its continuous extension. Using the center
\(\eta_j=\{\mu_P(j)-(1-\mu_P(j))\}/2\), we conclude
\[
\begin{aligned}
\operatorname{Var}_\star(D_j^\star)
&\le\mathbb E_\star[(D_j^\star-\eta_j)^2]\\
&\le\frac12\mathbb E_\star[(F_{j,+}-\mu_P(j))^2]
+\frac12\mathbb E_\star[(F_{j,-}-(1-\mu_P(j)))^2]\\
&\le\frac4{1+\zeta_j}.
\end{aligned}
\]
Together with the small-intensity case, this holds for every cell.

The second and fourth streams are independent. Their first and second moments consequently give
\[
\begin{aligned}
\operatorname{Var}_\star(V_j^\star D_j^\star)
&=\left(w_j^2+\frac{w_j}{m}\right)
\operatorname{Var}_\star(D_j^\star)
+\frac{w_j}{m}\{\mathbb E_\star[D_j^\star]\}^2\\
&\le\frac{4w_j^2}{1+\zeta_j}
+\frac{17w_j}{4m},
\end{aligned}
\]
where \(\{\mathbb E_\star[D_j^\star]\}^2\le1/4\).
Also,
\[
\frac{w_j^2}{1+mt_j}
\le\frac{2\delta}{m}w_j,
\qquad
\sum_{j\in\mathcal J}\frac{w_j^2}{1+mt_j}
\le\frac{2\delta^2}{m}.
\]
Since \(\delta+2\delta^2\le1\),
\[
\sum_{j\in\mathcal J}
\left(\frac{w_j}{m}+\frac{w_j^2}{1+mt_j}\right)
\le\frac1m.
\]
Replacing the coefficient \(4\) of the nonnegative quadratic term by \(17/4\) therefore yields
\[
\sum_{j\in\mathcal J}
\operatorname{Var}_\star(V_j^\star D_j^\star)
\le\frac{17}{4m}
=\frac{34}{n}.
\]
% lean: CausalSmith.Stat.MarNearcompleteFrontier.upper_streamVD_variance_sum_le

\item Define the total switching second moment by
\[
\mathcal V_{\mathrm{switch}}
=
\sum_{j\in\mathcal J}
\mathbb E_\star[
\{V_j^\star(G_j^\star-D_j^\star)\}^2].
\]
The identity
\[
G_j^\star-D_j^\star
=\mathbf1\{C_j^{\prime\star}\le B/4\}(H_j^\star-D_j^\star)
\]
and independence of the second, third, and fourth streams imply
\[
\mathbb E_\star[
\{V_j^\star(G_j^\star-D_j^\star)\}^2]
=
\left(w_j^2+\frac{w_j}{m}\right)\wp_j
\mathbb E_\star[(H_j^\star-D_j^\star)^2].
\]
Since \((u-v)^2\le2u^2+2v^2\) and
\(\mathbb E_\star[(D_j^\star)^2]\le1/4\),
\[
\mathbb E_\star[
\{V_j^\star(G_j^\star-D_j^\star)\}^2]
\le
\left(w_j^2+\frac{w_j}{m}\right)
\left\{2\wp_j\mathbb E_\star[(H_j^\star)^2]+\frac{\wp_j}{2}\right\}.
\]

For \(\zeta_j\le B\), use \(\wp_j\le1\), the light second-moment bound, and
\(w_j\le2\delta B/m\). Since \(e^{\ell_n/8}\ge1\), this gives
\[
\mathbb E_\star[
\{V_j^\star(G_j^\star-D_j^\star)\}^2]
\le
3e^{\ell_n/8}
\frac{4\delta^2B^2+2\delta B}{m^2}.
\]
Summing over at most \(4d\) cells,
\[
\sum_{\substack{j\in\mathcal J\\\zeta_j\le B}}
\mathbb E_\star[
\{V_j^\star(G_j^\star-D_j^\star)\}^2]
\le
12d\,e^{\ell_n/8}
\frac{4\delta^2B^2+2\delta B}{m^2}.
\]

For \(\zeta_j>B\), the pilot-weighted polynomial bound and
\(\wp_j\le e^{-\zeta_j/4}\le e^{-\zeta_j/8}\) give
\[
\mathbb E_\star[
\{V_j^\star(G_j^\star-D_j^\star)\}^2]
\le
\frac52\left(w_j^2+\frac{w_j}{m}\right)e^{-B/8}.
\]
Nonnegativity and \(\sum_jw_j\le\delta\) imply
\(\sum_jw_j^2\le\delta^2\). Hence
\[
\sum_{\substack{j\in\mathcal J\\\zeta_j>B}}
\mathbb E_\star[
\{V_j^\star(G_j^\star-D_j^\star)\}^2]
\le\frac52\left(\delta^2+\frac{\delta}{m}\right)e^{-B/8}
\le\frac{25}{n}.
\]
For the last inequality, use \(n\ge1\), \(\delta\le1\),
\(\delta/m\le8\), and
\[
e^{-B/8}=e^{-32\ell_n}\le e^{-\ell_n}
=\frac1{e+n}\le\frac1n.
\]
Thus
\[
\mathcal V_{\mathrm{switch}}
\le
12d\,e^{\ell_n/8}
\frac{4\delta^2B^2+2\delta B}{m^2}
+\frac{25}{n}.
\]

To convert the first term into the benchmark, define
\[
\Lambda_n=\frac{e^{\ell_n/8}\ell_n^3}{n}.
\]
Since \(\delta^2\le\delta\), \(\ell_n\le\ell_n^2\), and \(B=256\ell_n\),
\[
4\delta^2B^2+2\delta B
\le262656\,\delta\ell_n^2.
\]
Substitution of \(m=n/8\) gives the calibrated bound
\[
12d\,e^{\ell_n/8}
\frac{4\delta^2B^2+2\delta B}{m^2}
\le
C_{\mathrm{switch}}g_{n,d,q}\Lambda_n,
\qquad
C_{\mathrm{switch}}=12\cdot64\cdot262656.
\]
The sixth term of the exponential series yields
\[
\frac{(3\ell_n/4)^6}{720}\le e^{3\ell_n/4},
\qquad
\ell_n^6\le720(4/3)^6e^{3\ell_n/4}.
\]
Since \(e^{\ell_n}=e+n\le(e+1)n\),
\[
\Lambda_n^2
=\frac{e^{\ell_n/4}\ell_n^6}{n^2}
\le\frac{720(4/3)^6e^{\ell_n}}{n^2}
\le\frac{C_{\mathrm{log}}}{n}.
\]
Finally, the nonnegative square
\((g_{n,d,q}-C_{\mathrm{switch}}\Lambda_n)^2\) gives
\[
C_{\mathrm{switch}}g_{n,d,q}\Lambda_n
\le\frac{g_{n,d,q}^2}{2}
+\frac{C_{\mathrm{switch}}^2\Lambda_n^2}{2}
\le\frac{g_{n,d,q}^2}{2}
+\frac{C_{\mathrm{switch}}^2C_{\mathrm{log}}}{2n}.
\]
Consequently,
\[
\mathcal V_{\mathrm{switch}}
\le
\frac{g_{n,d,q}^2}{2}
+\frac{C_{\mathrm{switch}}^2C_{\mathrm{log}}/2+25}{n}.
\]
% lean: CausalSmith.Stat.MarNearcompleteFrontier.upper_stream_correction_second_moment_sum_rate

\item The decomposition
\[
V_j^\star G_j^\star
=V_j^\star D_j^\star
+V_j^\star(G_j^\star-D_j^\star)
\]
and the variance sum inequality give
\[
\operatorname{Var}_\star(V_j^\star G_j^\star)
\le
2\operatorname{Var}_\star(V_j^\star D_j^\star)
+2\mathbb E_\star[
\{V_j^\star(G_j^\star-D_j^\star)\}^2].
\]
Summing and using the two preceding bounds,
\[
\begin{aligned}
\operatorname{Var}_\star(W_{\mathrm{corr}})
&\le\frac{68}{n}+2\mathcal V_{\mathrm{switch}}\\
&\le
g_{n,d,q}^2+
\frac{C_{\mathrm{switch}}^2C_{\mathrm{log}}+118}{n}\\
&\le
\bigl(C_{\mathrm{switch}}^2C_{\mathrm{log}}+119\bigr)
r_{n,d,q}.
\end{aligned}
\]
Substitution into the aggregate risk inequality yields
\[
\begin{aligned}
\mathbb E_P[(\widehat\tau_{\mathrm{MM}}-\tau(P))^2]
&\le
\left\{
8\bigl(C_{\mathrm{switch}}^2C_{\mathrm{log}}+119\bigr)
+2C_{\mathrm{bias}}^2+104
\right\}r_{n,d,q}\\
&=C_{\mathrm{non}}r_{n,d,q}
\le C_Ur_{n,d,q}.
\end{aligned}
\]
Together with the complete-arrival and fallback branches, this establishes the asserted uniform bound for every admissible \(n,d,q,P\) and the stated independent sample.
% lean: CausalSmith.Stat.MarNearcompleteFrontier.upper_risk
\end{enumerate}
\end{proof}
